\documentclass[11pt,a4paper]{article}

%============================================================
%  Encodage, langue, polices
%============================================================
\usepackage[T1]{fontenc}
\usepackage[utf8]{inputenc}
\usepackage[main=french,provide=*]{babel}
\usepackage[expansion=false]{microtype}

%============================================================
%  Mathématiques
%============================================================
\usepackage{amsmath,amssymb,amsthm}
\usepackage{mathrsfs}   % \mathscr
\usepackage{stmaryrd}   % \llbracket ... \rrbracket
\usepackage{mathtools}

%============================================================
%  Mise en page
%============================================================
\usepackage[a4paper,margin=2.5cm]{geometry}
\usepackage{enumitem}
\usepackage{hyperref}
\hypersetup{
  colorlinks=true,
  linkcolor=black,
  urlcolor=blue!60!black,
  citecolor=black,
  pdftitle={Systèmes orthonormaux, séparabilité et mesure de comptage},
  pdfauthor={Stefan}
}

%============================================================
%  Environnements de théorèmes (style français)
%============================================================
\theoremstyle{plain}
\newtheorem{theoreme}{Théorème}[section]
\newtheorem{proposition}[theoreme]{Proposition}
\newtheorem{lemme}[theoreme]{Lemme}
\newtheorem{corollaire}[theoreme]{Corollaire}

\theoremstyle{definition}
\newtheorem{definition}[theoreme]{Définition}
\newtheorem{remarque}[theoreme]{Remarque}
\newtheorem{exemple}[theoreme]{Exemple}

\renewcommand{\proofname}{Démonstration}

%============================================================
%  Macros
%============================================================
\newcommand{\R}{\mathbb{R}}
\newcommand{\C}{\mathbb{C}}
\newcommand{\N}{\mathbb{N}}
\newcommand{\Q}{\mathbb{Q}}
\newcommand{\HH}{H}
\newcommand{\norm}[1]{\left\lVert #1 \right\rVert}
\newcommand{\abs}[1]{\left\lvert #1 \right\rvert}
\newcommand{\scal}[2]{\left\langle #1 , #2 \right\rangle}
\newcommand{\card}{\operatorname{card}}
\newcommand{\supp}{\operatorname{supp}}
\renewcommand{\Re}{\operatorname{Re}}
\newcommand{\muc}{\mu_{\mathrm{c}}}
\newcommand{\dimhilb}{\dim_{\mathrm{hilb}}}
\DeclareMathOperator{\Vect}{Vect}

\begin{document}

%============================================================
\begin{titlepage}
\centering
\vspace*{3cm}
{\Huge\bfseries Systèmes orthonormaux, séparabilité\\[0.4em] et mesure de comptage\par}
\vspace{1.5cm}
{\large Un système orthonormal d'un espace de Hilbert est-il dénombrable ?\par}
\vspace{2.5cm}
{\large\itshape Note de synthèse\par}
\vfill
{\large \today\par}
\end{titlepage}

\tableofcontents
\bigskip

%============================================================
\begin{abstract}
\noindent
On examine la question : \emph{un système orthonormal d'un espace de Hilbert
est-il nécessairement dénombrable ?} La réponse est négative en général, mais
affirmative sous l'hypothèse de \textbf{séparabilité}. On isole le mécanisme
purement métrique de la preuve — la séparation uniforme des vecteurs à
distance~$\sqrt{2}$ —, on montre que ce même mécanisme, \emph{lu à l'envers},
force la non-séparabilité de $\ell^2(I)$ lorsque $I$ est non dénombrable, puis on
identifie $\ell^2(I)$ comme l'espace $L^2$ associé à la \textbf{mesure de
comptage}. On clôt par une analyse du rôle exact joué par cette mesure, en
corrigeant au passage un point subtil sur les fonctions étagées.
\end{abstract}

%============================================================
\section{Le problème et son énoncé correct}

La proposition « tout système orthonormal d'un espace de Hilbert est
dénombrable » est \textbf{fausse} en toute généralité. L'énoncé correct fait
intervenir une hypothèse de séparabilité, qui porte toute la charge de la preuve.

\begin{definition}[Séparabilité]
Un espace métrique $\HH$ est \emph{séparable} s'il admet une partie dénombrable
dense, c'est-à-dire une partie $D \subset \HH$ au plus dénombrable telle que
$\overline{D} = \HH$.
\end{definition}

\begin{definition}[Système orthonormal]
Une famille $(e_i)_{i\in I}$ d'un espace de Hilbert $\HH$ est un \emph{système
orthonormal} si $\scal{e_i}{e_j} = \delta_{ij}$ pour tous $i,j \in I$.
\end{definition}

L'objectif de la première partie est de démontrer le résultat suivant, puis d'en
établir la nécessité de l'hypothèse.

\begin{theoreme}[Dénombrabilité sous séparabilité]\label{th:principal}
Dans un espace de Hilbert \textbf{séparable}, tout système orthonormal est
au plus dénombrable.
\end{theoreme}

%============================================================
\section{Le lemme de séparation et le théorème principal}

Le cœur de l'affaire est un calcul de distance élémentaire, valable dans tout
espace préhilbertien.

\begin{lemme}[Séparation uniforme à distance $\sqrt2$]\label{lem:sep}
Soit $(e_i)_{i\in I}$ un système orthonormal. Alors, pour $i \neq j$,
\[
  \norm{e_i - e_j} = \sqrt{2}.
\]
\end{lemme}

\begin{proof}
Par bilinéarité (ou sesquilinéarité) du produit scalaire et orthonormalité,
\[
  \norm{e_i - e_j}^2
   = \scal{e_i - e_j}{e_i - e_j}
   = \norm{e_i}^2 + \norm{e_j}^2 - 2\,\Re\scal{e_i}{e_j}
   = 1 + 1 - 0 = 2. \qedhere
\]
\end{proof}

La séparation uniforme se traduit géométriquement par une famille de boules deux
à deux disjointes, dont chacune isole un vecteur du système.

\begin{lemme}[Boules disjointes]\label{lem:boules}
Sous les hypothèses du lemme~\ref{lem:sep}, les boules ouvertes
$B_i = B\!\left(e_i, \tfrac{\sqrt2}{2}\right)$, $i \in I$, sont deux à deux
disjointes.
\end{lemme}

\begin{proof}
Supposons $x \in B_i \cap B_j$ avec $i \neq j$. L'inégalité triangulaire et le
lemme~\ref{lem:sep} donnent
\[
  \sqrt{2} = \norm{e_i - e_j}
   \le \norm{e_i - x} + \norm{x - e_j}
   < \frac{\sqrt2}{2} + \frac{\sqrt2}{2} = \sqrt2,
\]
ce qui est absurde. Donc $B_i \cap B_j = \varnothing$.
\end{proof}

\begin{proof}[Démonstration du théorème~\ref{th:principal}]
Soit $(e_i)_{i\in I}$ un système orthonormal de $\HH$ séparable, et soit
$D \subset \HH$ une partie dénombrable dense. Chaque boule $B_i$ du
lemme~\ref{lem:boules} est ouverte et non vide, donc rencontre $D$ par densité :
on peut choisir $d_i \in D \cap B_i$. Comme les $B_i$ sont deux à deux disjointes,
l'application
\[
  I \longrightarrow D, \qquad i \longmapsto d_i
\]
est \textbf{injective}. Il existe donc une injection de $I$ dans l'ensemble
dénombrable $D$, ce qui prouve que $I$ est au plus dénombrable.
\end{proof}

\begin{remarque}[La morale métrique]
L'orthonormalité n'intervient qu'\emph{une seule fois}, pour produire la
constante de séparation $\sqrt2$ (lemme~\ref{lem:sep}). Le reste est un fait
purement métrique :
\begin{quote}
\itshape
Dans un espace métrique séparable, toute partie $\delta$-séparée
(points deux à deux à distance $\ge \delta > 0$) est au plus dénombrable.
\end{quote}
On loge en effet une famille de boules disjointes de rayon $\delta/2$, chacune
capturant un point distinct de la partie dense.
\end{remarque}

%============================================================
\section{Nécessité de la séparabilité : la construction \texorpdfstring{$\ell^2(I)$}{l2(I)}}

Dire que l'hypothèse de séparabilité est \emph{nécessaire}, c'est nier
l'implication privée de son hypothèse : il faut exhiber un espace de Hilbert
\textbf{non séparable} portant un système orthonormal \textbf{non dénombrable}.
Le fait remarquable est que \emph{c'est le même calcul de séparation~$\sqrt2$}
qui, lu à l'endroit, donnait la dénombrabilité, et qui, lu à l'envers, produit la
non-séparabilité : les deux phénomènes sont les deux faces d'une même pièce.

\subsection{L'espace}

\begin{definition}[$\ell^2(I)$]
Pour un ensemble d'indices $I$ arbitraire, on pose
\[
  \ell^2(I) = \Bigl\{\, x = (x_i)_{i\in I} \in \C^{I} \ :\
    \sum_{i\in I} \abs{x_i}^2 < +\infty \,\Bigr\},
  \qquad
  \sum_{i\in I} \abs{x_i}^2 := \sup_{\substack{F \subset I \\ F \text{ fini}}}
    \sum_{i\in F} \abs{x_i}^2,
\]
muni du produit scalaire $\scal{x}{y} = \sum_{i\in I} x_i \overline{y_i}$.
C'est un espace de Hilbert (la complétude se démontre comme pour
$\ell^2(\N)$).
\end{definition}

\begin{remarque}[Chaque vecteur a un support dénombrable]\label{rem:support}
La finitude de $\sum_i \abs{x_i}^2$ force le support
$\{i : x_i \neq 0\}$ à être au plus dénombrable. En effet, pour $n \ge 1$
l'ensemble $\{i : \abs{x_i} > 1/n\}$ est fini (sinon la somme divergerait), et
\[
  \{i : x_i \neq 0\} = \bigcup_{n \ge 1} \Bigl\{ i : \abs{x_i} > \tfrac1n \Bigr\}
\]
est réunion dénombrable d'ensembles finis. \emph{Chaque vecteur individuel ne
« voit » qu'un nombre dénombrable de coordonnées.} On pourrait croire l'espace
« moralement séparable » — c'est faux, et toute la finesse est là : aucun jeu
\emph{dénombrable} de vecteurs ne les voit \emph{tous} à la fois.
\end{remarque}

\subsection{Le système orthonormal non dénombrable}

Pour $i \in I$, soit $e_i = (\delta_{ij})_{j\in I}$ le vecteur valant $1$ en
position $i$ et $0$ ailleurs. Chaque $e_i \in \ell^2(I)$ (support réduit à un
point) et $\scal{e_i}{e_j} = \delta_{ij}$ : la famille $(e_i)_{i\in I}$ est
orthonormale, indexée par $I$.

\begin{theoreme}[Non-séparabilité]\label{th:nonsep}
Si $I$ est non dénombrable, alors $\ell^2(I)$ possède un système orthonormal non
dénombrable et n'est pas séparable.
\end{theoreme}

\begin{proof}
La famille $(e_i)_{i\in I}$ est orthonormale de cardinal $\card(I)$ non
dénombrable, ce qui règle le premier point. Par le lemme~\ref{lem:sep} elle est
$\sqrt2$-séparée, donc les boules $B_i = B\!\left(e_i, \tfrac{\sqrt2}{2}\right)$
sont deux à deux disjointes (lemme~\ref{lem:boules}). Supposons par l'absurde
$\ell^2(I)$ séparable, de partie dense dénombrable $D$. Par densité, chaque $B_i$
contient un $d_i \in D$ ; par disjonction des boules, $i \mapsto d_i$ est une
injection de $I$ dans $D$. On obtiendrait une injection de $I$ (non dénombrable)
dans $D$ (dénombrable) : contradiction.
\end{proof}

\begin{remarque}[La même mécanique, à rôles échangés]
La preuve du théorème~\ref{th:nonsep} est \emph{mot pour mot} celle du
théorème~\ref{th:principal}, à ceci près qu'on a échangé les rôles de
l'hypothèse et de la conclusion.
\begin{itemize}[nosep,leftmargin=1.4em]
  \item \textbf{Sens direct} : on \emph{dispose} d'un $D$ dénombrable, et on en
        déduit que le système orthonormal, injecté dans $D$, est dénombrable.
  \item \textbf{Sens réciproque} : on \emph{dispose} d'un système orthonormal non
        dénombrable, on l'injecte dans tout candidat $D$ à la densité, et on en
        déduit que $D$ ne peut être dénombrable.
\end{itemize}
La famille $\sqrt2$-séparée est l'obstruction commune : tantôt objet contraint
par la petitesse de l'espace, tantôt témoin de sa grosseur.
\end{remarque}

\begin{remarque}[Un exemple universel, non ad hoc]
$\ell^2(I)$ n'est pas un exemple artificiel. Tout espace de Hilbert admet une
base hilbertienne, et deux bases hilbertiennes d'un même espace ont même cardinal
(la \emph{dimension hilbertienne} $\dimhilb$). Si $\dimhilb \HH = \kappa$, alors
$\HH \cong \ell^2(I)$ isométriquement avec $\card(I) = \kappa$. La dichotomie est
donc totale :
\[
  \HH \text{ séparable}
  \iff \dimhilb \HH \le \aleph_0
  \iff \HH \cong \C^{n} \ \text{ou}\ \ell^2(\N).
\]
Les systèmes orthonormaux non dénombrables vivent très exactement — et nulle part
ailleurs — dans les $\ell^2(I)$ de dimension hilbertienne non dénombrable.
\end{remarque}

%============================================================
\section{\texorpdfstring{$\ell^2(I)$}{l2(I)} comme espace \texorpdfstring{$L^2$}{L2} de la mesure de comptage}

La « somme sur un ensemble d'indices arbitraire », en apparence \emph{ad hoc},
est en réalité un cas particulier de la théorie de l'intégration.

\begin{definition}[Mesure de comptage]
Sur l'espace mesurable $(I, \mathscr{P}(I))$, la \emph{mesure de comptage}
$\muc$ est définie par $\muc(A) = \card(A) \in [0,+\infty]$.
\end{definition}

\begin{proposition}[Identification]\label{prop:ident}
On a l'égalité isométrique
\[
  \ell^2(I) = L^2\bigl(I, \mathscr{P}(I), \muc\bigr),
\]
et pour toute $f \ge 0$ sur $I$,
\[
  \int_I f \, d\muc \;=\; \sum_{i\in I} f(i)
     \;:=\; \sup_{\substack{F \subset I \\ F \text{ fini}}} \sum_{i\in F} f(i).
\]
\end{proposition}

Cette dernière formule n'est pas une convention : c'est un \textbf{théorème}. Le
membre de gauche est l'intégrale de Lebesgue abstraite (borne supérieure des
intégrales des fonctions étagées $\le f$) ; comme les fonctions étagées
\emph{intégrables} sont ici à support fini (section~\ref{sec:etagees}), ce
supremum se calcule sur les sommes finies. Ainsi $\sum_i \abs{x_i}^2 =
\int_I \abs{x}^2 \, d\muc$ et $\scal{x}{y} = \int_I x\,\overline{y}\, d\muc$.

\subsection{Deux simplifications offertes par $\muc$}

\begin{itemize}[leftmargin=1.4em]
  \item \textbf{La tribu est le powerset entier.} Toute fonction $I \to \C$ est
        mesurable ; aucune vérification de mesurabilité n'est jamais requise.
  \item \textbf{Aucun quotient par les négligeables.} Un ensemble est
        $\muc$-négligeable si et seulement s'il est vide :
        $\muc(A) = 0 \iff A = \varnothing$. Il n'y a donc pas d'égalité
        « presque partout » non triviale, et $\ell^2(I)$ est un espace de
        \emph{véritables} fonctions, non de classes d'équivalence — contrairement
        à $L^2([0,1],\text{Leb})$. C'est pourquoi $e_i$ est \emph{le} vecteur
        valant $1$ en $i$, sans ambiguïté.
\end{itemize}

L'orthonormalité se relit alors de façon purement mesurée : $e_i = \mathbf{1}_{\{i\}}$ et
\[
  \scal{e_i}{e_j} = \int_I \mathbf{1}_{\{i\}}\mathbf{1}_{\{j\}}\, d\muc
     = \muc(\{i\}\cap\{j\}) = \delta_{ij}.
\]
L'orthonormalité est exactement le fait que chaque singleton a masse
$\muc(\{i\}) = 1$.

\subsection{La contrepartie : perte de la $\sigma$-finitude}

\begin{proposition}[Non-$\sigma$-finitude]
Si $I$ est non dénombrable, $\muc$ n'est pas $\sigma$-finie : les parties de
mesure finie sont les parties finies, et une réunion dénombrable de parties
finies reste dénombrable, donc $\neq I$.
\end{proposition}

Les théorèmes exigeant la $\sigma$-finitude (Radon--Nikodym sous forme standard,
Fubini--Tonelli, certains cas de la dualité $L^p$--$L^q$) tombent alors. En
revanche, un fait de $\sigma$-finitude \emph{locale} survit toujours, et il
explique la remarque~\ref{rem:support}.

\begin{proposition}[Support $\sigma$-fini universel des fonctions $L^2$]
Soit $(X,\mathscr{A},\mu)$ un espace mesuré \emph{quelconque} et $f \in L^2(\mu)$.
Alors $\{f \neq 0\}$ est $\sigma$-fini.
\end{proposition}

\begin{proof}
Par l'inégalité de Bienaymé--Tchebychev, pour tout $n \ge 1$,
\[
  \mu\bigl(\abs{f} > \tfrac1n\bigr)
   \le n^2 \int_X \abs{f}^2 \, d\mu < +\infty,
\]
donc $\{f \neq 0\} = \bigcup_{n\ge1} \{\abs{f} > \tfrac1n\}$ est réunion
dénombrable d'ensembles de mesure finie.
\end{proof}

Pour $\muc$, « $\sigma$-fini » signifie « réunion dénombrable de parties finies »,
c'est-à-dire « dénombrable ». Le support dénombrable de chaque $x \in \ell^2(I)$
n'est donc \emph{pas} une curiosité : c'est le support $\sigma$-fini universel des
fonctions $L^2$, spécialisé à la mesure de comptage. Ce qui est frappant, c'est
que la $\sigma$-finitude reste vraie \emph{fonction par fonction} tout en
échouant \emph{globalement} — et c'est exactement cet écart qui fabrique la
non-séparabilité (théorème~\ref{th:nonsep}).

\begin{remarque}
La structure hilbertienne, elle, ne souffre de rien : $L^2$ d'un espace mesuré
\emph{quelconque} est toujours complet (Riesz--Fischer n'exige pas la
$\sigma$-finitude). Seuls certains outils annexes de la théorie de la mesure
deviennent indisponibles.
\end{remarque}

%============================================================
\section{Fonctions étagées et effondrement de l'intégrale en somme}
\label{sec:etagees}

Un point mérite d'être énoncé avec soin, car sa formulation naïve est fausse.

\begin{remarque}[Un piège]
Une fonction étagée (mesurable, ne prenant qu'un nombre \emph{fini} de valeurs)
n'est \textbf{pas} nécessairement à support fini. Contre-exemple :
$\mathbf{1}_I$ sur $I$ infini est étagée — elle ne prend que la valeur $1$ — mais
son support est $I$ tout entier.
\end{remarque}

L'énoncé correct fait intervenir l'intégrabilité.

\begin{lemme}[Étagée intégrable $\Rightarrow$ support fini]\label{lem:etagee}
Soit $s$ une fonction étagée sur $(I, \mathscr{P}(I), \muc)$. Si
$\int_I \abs{s}\, d\muc < +\infty$, alors $s$ est à support fini. Plus
précisément,
\[
  \int_I \abs{s}\, d\muc \;\ge\; \Bigl(\min_{k} \abs{c_k}\Bigr)
     \cdot \card\bigl(\supp s\bigr),
\]
où les $c_k$ sont les valeurs non nulles de $s$.
\end{lemme}

\begin{proof}
Écrivons $s = \sum_{k=1}^{n} c_k \, \mathbf{1}_{A_k}$ avec les $c_k \neq 0$ deux à
deux distincts (on écarte le morceau nul) et $A_k = \{s = c_k\}$ disjoints. Comme
$\muc(\{i\}) = 1$,
\[
  \int_I \abs{s}\, d\muc = \sum_{k=1}^{n} \abs{c_k}\, \muc(A_k)
     = \sum_{k=1}^{n} \abs{c_k}\, \card(A_k).
\]
Chaque $\abs{c_k} > 0$, donc cette somme finie de termes positifs est finie si et
seulement si chaque $\card(A_k) < \infty$, c'est-à-dire chaque $A_k$ fini. Alors
$\supp s = \bigcup_{k=1}^{n} A_k$, réunion \emph{finie} d'ensembles finis, est
fini. La minoration annoncée en découle en majorant chaque $\abs{c_k}$ par le
minimum et en sommant les cardinaux. Contraposée : support infini $\Rightarrow$
l'un des $A_k$ (à valeur non nulle) est infini $\Rightarrow$ intégrale infinie.
\end{proof}

C'est ce lemme qui fait s'effondrer l'intégrale abstraite en somme non ordonnée.

\begin{proof}[Démonstration de la proposition~\ref{prop:ident} (formule intégrale)]
Soit $f \ge 0$. On montre l'identité
\[
  \int_I f\, d\muc = \sup_{F}\ \sum_{i\in F} f(i),
\]
le supremum portant sur les parties finies $F \subset I$.

\medskip
\noindent\textbf{Inégalité $\ge$.} Pour $F$ fini, $f\mathbf{1}_F$ est étagée à
support fini, $0 \le f\mathbf{1}_F \le f$, et
$\int_I f\mathbf{1}_F\, d\muc = \sum_{i\in F} f(i)$. Donc
$\int_I f\, d\muc \ge \sum_{i\in F} f(i)$, puis $\int_I f\, d\muc \ge \sup_F$.

\medskip
\noindent\textbf{Inégalité $\le$.} Soit $s$ étagée avec $0 \le s \le f$.
\begin{itemize}[nosep,leftmargin=1.4em]
  \item Si $\int_I s\, d\muc = +\infty$, alors $s$ a un support infini
        (lemme~\ref{lem:etagee}) sur lequel $f \ge s \ge \min_k c_k > 0$, d'où
        $\sup_F \sum_{i\in F} f(i) = +\infty$ : l'inégalité est triviale.
  \item Si $\int_I s\, d\muc < +\infty$, le lemme~\ref{lem:etagee} donne
        $s$ à support fini ; en posant $F = \supp s$,
        \[
          \int_I s\, d\muc = \sum_{i \in F} s(i)
             \le \sum_{i\in F} f(i) \le \sup_F .
        \]
\end{itemize}
En prenant le supremum sur les étagées $s \le f$, on obtient
$\int_I f\, d\muc \le \sup_F$. D'où l'égalité.
\end{proof}

\begin{remarque}[Le contraste qui isole le rôle de $\muc$]
Le support fini des étagées intégrables est un trait \emph{spécifique} de la
mesure de comptage. Pour la mesure de Lebesgue sur $\R$, une étagée intégrable a
un support de \emph{mesure} finie, ce qui n'a rien à voir avec un support fini en
\emph{cardinal} : $\mathbf{1}_{[0,1]}$ est intégrable, de support non dénombrable.
La spécificité de $\muc$ est l'égalité $\muc(A) = \card(A)$ : « intégrale finie »
s'y traduit \emph{mot pour mot} en « support fini », l'égalité entre masse et
cardinal transportant la finitude de l'une à l'autre. C'est exactement cette
traduction qui fait de $\ell^2(I)$ un espace de \emph{sommes} plutôt que
d'\emph{intégrales}.
\end{remarque}

%============================================================
\section{Synthèse}

Le fil conducteur est une tension unique, quantifiée par la distance $\sqrt2$,
entre deux exigences contradictoires :
\begin{itemize}[leftmargin=1.4em]
  \item la \textbf{séparabilité}, contrainte de petitesse topologique : tout point
        est approchable par une réserve dénombrable fixée ;
  \item un \textbf{système orthonormal}, configuration maximalement étalée : ses
        points sont mutuellement à distance $\sqrt2$, chacun occupant son propre
        territoire — une boule de rayon $\sqrt2/2$ que nul autre ne partage.
\end{itemize}
Un espace séparable ne loge qu'un nombre dénombrable de tels territoires
disjoints ; dès qu'il y en a un nombre non dénombrable, la réserve dense doit
elle-même être non dénombrable. Tout le reste — l'identification à
$L^2(I,\muc)$, la perte de $\sigma$-finitude, le support fini des étagées
intégrables — n'est qu'une lecture, à travers la théorie de la mesure, de ce même
fait métrique.

\begin{center}
\fbox{\parbox{0.92\linewidth}{\centering
\medskip
$\HH$ séparable
\quad$\Longleftrightarrow$\quad
tout système orthonormal de $\HH$ est dénombrable
\\[0.6em]
$\Longleftrightarrow\ \dimhilb \HH \le \aleph_0
   \ \Longleftrightarrow\ \HH \cong \C^{n}\ \text{ou}\ \ell^2(\N).$
\medskip
}}
\end{center}

\end{document}
